% Options for packages loaded elsewhere
\PassOptionsToPackage{unicode}{hyperref}
\PassOptionsToPackage{hyphens}{url}
\PassOptionsToPackage{dvipsnames,svgnames,x11names}{xcolor}
\documentclass[
  11pt]{article}
\usepackage{xcolor}
\usepackage[margin=1in]{geometry}
\usepackage{amsmath,amssymb}
\setcounter{secnumdepth}{-\maxdimen} % remove section numbering
\usepackage{iftex}
\ifPDFTeX
  \usepackage[T1]{fontenc}
  \usepackage[utf8]{inputenc}
  \usepackage{textcomp} % provide euro and other symbols
\else % if luatex or xetex
  \usepackage{unicode-math} % this also loads fontspec
  \defaultfontfeatures{Scale=MatchLowercase}
  \defaultfontfeatures[\rmfamily]{Ligatures=TeX,Scale=1}
\fi
\usepackage{lmodern}
\ifPDFTeX\else
  % xetex/luatex font selection
\fi
% Use upquote if available, for straight quotes in verbatim environments
\IfFileExists{upquote.sty}{\usepackage{upquote}}{}
\IfFileExists{microtype.sty}{% use microtype if available
  \usepackage[]{microtype}
  \UseMicrotypeSet[protrusion]{basicmath} % disable protrusion for tt fonts
}{}
\makeatletter
\@ifundefined{KOMAClassName}{% if non-KOMA class
  \IfFileExists{parskip.sty}{%
    \usepackage{parskip}
  }{% else
    \setlength{\parindent}{0pt}
    \setlength{\parskip}{6pt plus 2pt minus 1pt}}
}{% if KOMA class
  \KOMAoptions{parskip=half}}
\makeatother
\usepackage{longtable,booktabs,array}
\usepackage{caption}
\captionsetup[table]{skip=6pt}
\newcounter{none} % for unnumbered tables
\usepackage{calc} % for calculating minipage widths
% Correct order of tables after \paragraph or \subparagraph
\usepackage{etoolbox}
\makeatletter
\patchcmd\longtable{\par}{\if@noskipsec\mbox{}\fi\par}{}{}
\makeatother
\setlength{\emergencystretch}{3em} % prevent overfull lines
\providecommand{\tightlist}{%
  \setlength{\itemsep}{0pt}\setlength{\parskip}{0pt}}
\usepackage{xurl}
\usepackage{bookmark}
\IfFileExists{xurl.sty}{\usepackage{xurl}}{} % add URL line breaks if available
\urlstyle{same}
% fallback for those not using the hyperref driver hyperxmp:
\makeatletter
\@ifundefined{xmpquote}{\newcommand{\xmpquote}[1]{#1}}{}
\makeatother
\hypersetup{
  pdftitle={Stolen keys are worthless: placeholder credentials that become real only inside encryption{,} with tripwires},
  pdfauthor={Spencer Nunamaker},
  colorlinks=true,
  linkcolor={black},
  filecolor={Maroon},
  citecolor={Blue},
  urlcolor={blue},
  pdfcreator={LaTeX via pandoc}}

\title{Stolen keys are worthless:
placeholder credentials that become real
only inside encryption, with tripwires}
\author{Spencer Nunamaker\\
{junovhs@gmail.com ·
\url{https://github.com/junovhs/keyhandoff}}}
\date{}

\begin{document}
\maketitle

\section*{Abstract}\label{abstract}
\addcontentsline{toc}{section}{Abstract}

Developers keep API keys where their
tools can read them: in \texttt{.env}
files, in environment variables, in
shell history and logs. Supply-chain
malware and, increasingly, AI coding
agents read the same places. Vaults and
run wrappers (\texttt{op\ run},
\texttt{doppler\ run}) keep keys out of
files but hand the real key to the
program, where anything in its process
can take it. Placeholder proxies keep
the real key out of the program but
require routing traffic through a
man-in-the-middle and only speak HTTP.

We present a design in which a program
never holds the key it uses. It holds a
placeholder with the key's prefix and
length. A hook in the program's own TLS
library swaps the real key in at the
moment a write is encrypted, and only on
a connection whose certificate chain the
hook has verified itself, against the
system's roots, as one of the key's
servers. Because nothing legitimate
sends a placeholder anywhere else, every
placeholder is also a tripwire. A write
carrying one elsewhere, as written or in
common encodings, is refused before a
byte leaves. The run that tried is
named, the destination is refused from
then on, and a signed warning reaches
teammates through a folder they already
share. Each run gets fresh placeholders,
so one found later traces back to its
run. We implement the design in
KeyHandoff, a small desktop key manager.

We evaluate coverage across 15 client
and protocol combinations, 11
reconstructed attacks drawn from
documented supply-chain campaigns and
AI-agent misuse, overhead, and an
adversarial review.

\begin{itemize}
\tightlist
\item
  \textbf{Coverage:} 12 of the 15
  combinations are protected, including
  HTTP/2, where we rewrite
  HPACK-compressed headers {[}17{]}. The
  other three never release the key.
\item
  \textbf{Attacks:} in every
  reconstructed attack the attacking
  code held only a placeholder, and no
  real key left the machine. Every send
  was stopped and attributed.
\item
  \textbf{Overhead:} a swapped request
  costs about 0.26 ms; anything else
  costs about 2 µs.
\item
  \textbf{Adversarial review:} it found,
  and we fixed, a way for in-process
  code to obtain the key through the
  hook's own token. It also states
  precisely what still does: code that
  targets the hook in memory, and use of
  the key at its own server.
\end{itemize}

\section*{1. The
problem}\label{the-problem}
\addcontentsline{toc}{section}{1. The
problem}

A key is safe until a program needs it.
At that point it ends up somewhere other
code can read:

\begin{itemize}
\tightlist
\item
  \textbf{Files and the environment.}
  \texttt{.env} files, shell profiles
  and CI variables hold keys in plain
  text. Stealers target exactly these.
  The August 2025 compromise of
  \texttt{nx} packages inventoried
  environment variables and
  \texttt{.env} files, asked locally
  installed AI command-line tools to
  search for more, and uploaded the haul
  triple-base64-encoded to a public
  GitHub repository created in the
  victim's own account {[}1, 2{]}. The
  September 2025 Shai-Hulud worm, run
  from npm install scripts, harvested
  environment variables and cloud keys,
  published them double-base64-encoded
  to a public repository named
  \texttt{Shai-Hulud}, and planted a
  GitHub Action that sent each
  repository's secrets to a webhook
  {[}3{]}.
\item
  \textbf{AI coding agents.} An agent
  with shell access can print, encode or
  upload a key. It may do so because it
  was asked, or because it was told to
  by text it read: an issue, a web page,
  a README. An agent can also run the
  wrapper itself
  (\texttt{op\ run\ -\/-\ sh\ -c\ \textquotesingle{}env\textquotesingle{}}).
\item
  \textbf{Logs and transcripts.} Keys
  printed while debugging end up in
  terminal scrollback, CI logs and agent
  transcripts.
\end{itemize}

Existing tools each remove part of this
exposure:

{\def\LTcaptype{none} % do not increment counter
\begin{longtable}[]{@{}
  >{\raggedright\arraybackslash}p{(\linewidth - 8\tabcolsep) * \real{0.2000}}
  >{\raggedright\arraybackslash}p{(\linewidth - 8\tabcolsep) * \real{0.2000}}
  >{\raggedright\arraybackslash}p{(\linewidth - 8\tabcolsep) * \real{0.2000}}
  >{\raggedright\arraybackslash}p{(\linewidth - 8\tabcolsep) * \real{0.2000}}
  >{\raggedright\arraybackslash}p{(\linewidth - 8\tabcolsep) * \real{0.2000}}@{}}
\toprule\noalign{}
\begin{minipage}[b]{\linewidth}\raggedright
Approach
\end{minipage} &
\begin{minipage}[b]{\linewidth}\raggedright
Key in a file
\end{minipage} &
\begin{minipage}[b]{\linewidth}\raggedright
Key in the program
\end{minipage} &
\begin{minipage}[b]{\linewidth}\raggedright
Traffic
\end{minipage} &
\begin{minipage}[b]{\linewidth}\raggedright
Protocols
\end{minipage} \\
\midrule\noalign{}
\endhead
\bottomrule\noalign{}
\endlastfoot
\texttt{.env} & yes & yes & direct &
any \\
Vault + run wrapper (\texttt{op\ run},
\texttt{doppler\ run},
\texttt{infisical\ run}) {[}4--6{]} & no
& yes & direct & any \\
Placeholder proxy (agent sandboxes,
\texttt{agent-vault},
\texttt{iron-proxy}, fnox, GhostKey)
{[}7--11{]} & no & no & through a MITM
proxy and its CA & HTTP \\
Kernel uprobe swap (Kloak) {[}14{]} & no
& no & direct & TLS, in a cluster, with
privileges \\
\textbf{This work} & no & no (except
during the swap) & direct & anything
over TLS the hook covers \\
\end{longtable}
}

\section*{2. Threat
model}\label{threat-model}
\addcontentsline{toc}{section}{2. Threat
model}

We defend a developer's own computer,
which is not compromised at the
administrator level. The attacker runs
code as the developer, either inside a
program the developer starts (a
malicious dependency, an agent's
command) or beside it. That code can
read the program's environment, files
and output, and can make network
connections. It cannot change the
kernel, read other users' memory, or
break TLS.

The goal is that the key never leaves
the machine except to its own server,
inside a TLS connection whose
certificate the hook verified.
Everything the attacker can take is a
placeholder. Placeholders die when their
run ends, and every attempt to send one
elsewhere is visible.

Out of scope, and said plainly in
\texttt{docs/threats.md}:

\begin{itemize}
\tightlist
\item
  \textbf{Use at the key's own server.}
  Code in the program can make requests
  to the key's own server as the
  developer, exactly as any program
  allowed the key could.
\item
  \textbf{Code that targets the hook in
  memory.} The key is in the program's
  memory for the duration of each swap.
\item
  \textbf{Administrator-level
  compromise.}
\end{itemize}

\section*{3. Design}\label{design}
\addcontentsline{toc}{section}{3.
Design}

The design has four parts. A \textbf{key
service}, a long-running process of the
developer's, holds the keys and decides
which programs may use them. A
\textbf{runner} starts a program with
its keys. A \textbf{hook}, loaded into
that program, swaps keys in and watches
for placeholders leaving. \textbf{Team
folders}, ordinary folders teammates
already share, carry signed warnings
between their key services. In our
prototype, KeyHandoff is the key
service, and its
\texttt{keyhandoff\ run} command is the
runner.

\subsection*{3.1 Running a program with
keys}\label{running-a-program-with-keys}
\addcontentsline{toc}{subsection}{3.1
Running a program with keys}

The runner reads pointers to keys (in
our prototype, lines like
\texttt{API\_KEY=handoff:Name}) from
\texttt{.env} and the environment. It
asks the key service for those keys over
a per-user local socket, naming the
exact command it will start, then starts
the command.

\textbf{Approval is bound to the
command.} An approval covers the
program's path, its arguments, and the
SHA-256 of every file the arguments
name. An edited script, or the same
interpreter running other code
(\texttt{python\ -c\ …}), has to ask
again. The runner's claim about the
command is believed only when the
request comes from the runner's own
executable, which the key service checks
by file identity or content.

\textbf{Prompts are AI-aware.} The key
service walks the requesting process's
parent chain. If Claude Code, Cursor,
Codex, Copilot, Aider, Windsurf or
Gemini CLI appears in it, saved
approvals don't count: the prompt names
the agent and the exact command, and
offers ``Allow for 30 minutes'', which
is never saved, instead of ``Always''.
Name-based detection is a warning; the
binding to the command is the guarantee.

\textbf{Output is masked.} The program's
stdout and stderr, and those of
everything it starts, pass through a
masker. It hides the key or placeholder
printed plain, base64 (at each byte
offset), hex, percent-encoded or
reversed, and any run of 12 or more of
those characters, across chunk
boundaries.

\subsection*{3.2
Placeholders}\label{placeholders}
\addcontentsline{toc}{subsection}{3.2
Placeholders}

Each key records the servers it is for.
They are guessed from well-known
prefixes (\texttt{sk\_live\_} means
\texttt{api.stripe.com},
\texttt{sk-ant-} means
\texttt{api.anthropic.com}) and can be
edited. For such a key, the program gets
a placeholder instead of the key: the
same prefix and length, with random
characters of the same classes after it,
fresh for every run. Same length means
client libraries accept it and
\texttt{Content-Length} doesn't change.
A key its provider signs with rather
than sends (AWS) gets no default server,
because a placeholder cannot work for
it.

\subsection*{3.3 The
swap}\label{the-swap}
\addcontentsline{toc}{subsection}{3.3
The swap}

On Linux the runner preloads a small C
library into the program
(\texttt{LD\_PRELOAD}). It interposes
OpenSSL's \texttt{SSL\_write} and
\texttt{SSL\_write\_ex}, resolving
OpenSSL at run time, so the library
needs no headers and works with OpenSSL
1.1 and 3. When a write holds a
placeholder, the hook does four things:

\begin{enumerate}
\def\labelenumi{\arabic{enumi}.}
\tightlist
\item
  \textbf{Verify the chain itself.} It
  checks the peer's certificate chain
  against the system's default roots
  and, if set, the key's own CA. It
  never uses the program's trust
  settings, \texttt{SSL\_CERT\_FILE} or
  the verification result the program
  saw, because code in the program can
  change all three.
\item
  \textbf{Check the host.} It checks the
  leaf certificate against the key's
  hosts with \texttt{X509\_check\_host}.
  The name asked for in SNI and the
  \texttt{Host} header play no part.
\item
  \textbf{Fetch the key.} It asks the
  key service for the key, presenting
  the run's token. The key service
  answers only processes descended from
  the run's runner, and only while the
  runner is connected.
\item
  \textbf{Swap.} It writes the key into
  a private copy of the buffer, writes
  that copy (with
  \texttt{SSL\_MODE\_ACCEPT\_MOVING\_WRITE\_BUFFER},
  so a retried write can come from
  another buffer), and wipes it.
\end{enumerate}

The verification verdict is cached per
connection for a minute. The cache key
is the connection object together with
the leaf certificate's signature bytes,
so a reused object talking to another
server misses.

\textbf{HTTP/2.} HPACK {[}17{]}
compresses headers, usually with a
static Huffman code, before they reach
\texttt{SSL\_write}, so a placeholder in
an HTTP/2 header is never there as
written. On a connection that negotiated
\texttt{h2}, the hook follows the frames
from the connection preface. It decodes
each complete HEADERS block's string
literals, swaps the placeholder where
the connection is verified, and
re-encodes each string with the same
representation. Because the key is as
long as its placeholder, the dynamic
tables on both sides stay in step. The
Huffman table was not typed in. We
measured each symbol's code length from
the system's libnghttp2 encoder, rebuilt
the canonical code from those lengths,
and checked it byte for byte against
libnghttp2 on 1,898 random strings; the
generator is part of the artifact (§9).

\textbf{HTTP Basic auth}
(\texttt{curl\ -u\ key:}) sends the key
base64-encoded. At the key's own server,
the hook decodes the base64 run, swaps,
and encodes it again at the same length.

\textbf{Node.js} links its own OpenSSL,
out of the preloaded hook's reach. The
runner adds a \texttt{-\/-require}
module to \texttt{NODE\_OPTIONS}. It
wraps the TLS socket write path, waits
for the handshake, verifies the chain
itself (with
\texttt{crypto.X509Certificate}, against
Node's root store and the key's CA), and
swaps or refuses.

\textbf{The token.} The hook reads its
token at load and overwrites it in
place. Just unsetting it would leave it
in \texttt{/proc/self/environ} (§6.4).
It passes the token on to programs the
hook starts through \texttt{exec},
\texttt{posix\_spawn}, \texttt{system}
and \texttt{popen}.

\subsection*{3.4
Tripwires}\label{tripwires}
\addcontentsline{toc}{subsection}{3.4
Tripwires}

A placeholder headed anywhere but its
key's servers is a leak in the act.

\textbf{What is refused.} The hook
refuses, before any of it leaves:

\begin{itemize}
\tightlist
\item
  a TLS write holding a placeholder on a
  connection not verified for its key;
\item
  any write holding one on a plain
  internet socket (\texttt{send},
  \texttt{sendto}, \texttt{sendmsg},
  \texttt{write}, \texttt{writev});
\item
  a name lookup (\texttt{getaddrinfo})
  containing one.
\end{itemize}

It recognises the placeholder as
written, in standard base64 nested up to
three levels deep at every byte offset,
in URL-safe base64, hex,
percent-encoding and reversed. Those 47
patterns are scanned in one pass over a
hashed set of 4-byte windows, probed at
a stride of the shortest pattern length
minus three.

\textbf{What happens next.} The key
service keeps an alert naming the key,
the command that tried, and the
destination. It refuses that destination
for later runs: by name, by certificate
name, and by address. It also writes a
signed warning into each team folder the
developer is on. Teammates' key services
refuse destinations from warnings signed
by a member of the team.

\textbf{Lineage.} Every run's
placeholders are recorded, so a
placeholder found later says which key,
command and time it came from, even
base64-encoded up to three times. A
placeholder can turn up in a log, a
commit or an agent transcript.

\section*{4.
Implementation}\label{implementation}
\addcontentsline{toc}{section}{4.
Implementation}

KeyHandoff is a Tauri app with a Rust
core and React screens. It is both the
key service and, as
\texttt{keyhandoff\ run}, the runner.
The parts described here are:

\begin{itemize}
\tightlist
\item
  \textbf{Runner:} \texttt{run.rs}, and
  \texttt{mask.rs} for output masking.
\item
  \textbf{Placeholders and hosts:}
  \texttt{placeholder.rs}.
\item
  \textbf{The C hook:}
  \texttt{src-tauri/hook/keyhandoff\_tls.c},
  about 1,400 lines, with its HPACK
  Huffman table generated by
  \texttt{gen\_hpack\_huffman.py},
  compiled by \texttt{build.rs} and
  embedded in the executable.
\item
  \textbf{The Node hook:}
  \texttt{keyhandoff\_node.cjs}, about
  210 lines.
\item
  \textbf{Signed team warnings:} a new
  object kind in the sealing crate,
  \texttt{crates/seal/src/warning.rs}.
\end{itemize}

The hook is written out on first use to
a private folder under
\texttt{\$XDG\_RUNTIME\_DIR}, named by
its contents.

\section*{5.
Evaluation}\label{evaluation}
\addcontentsline{toc}{section}{5.
Evaluation}

Everything below is regenerated by
\texttt{eval/run.sh} and, for coverage
and attacks, by CI on every push. CI
runs on GitHub's Ubuntu 22.04 runner;
the Windows build and tests also run
there, but no Windows cell is claimed.
The numbers below are from a Ryzen 9
3900X running Pop!\_OS 22.04 with
OpenSSL 3.0.2.

\subsection*{5.1
Coverage}\label{coverage}
\addcontentsline{toc}{subsection}{5.1
Coverage}

The same keyed request goes to a server
whose certificate is valid for the key's
host, and to a stranger with a valid
certificate for another name. Both
servers use the same test CA, configured
as the key's CA.

{\def\LTcaptype{none} % do not increment counter
\begin{longtable}[]{@{}
  >{\raggedright\arraybackslash}p{(\linewidth - 8\tabcolsep) * \real{0.2000}}
  >{\raggedright\arraybackslash}p{(\linewidth - 8\tabcolsep) * \real{0.2000}}
  >{\raggedright\arraybackslash}p{(\linewidth - 8\tabcolsep) * \real{0.2000}}
  >{\raggedright\arraybackslash}p{(\linewidth - 8\tabcolsep) * \real{0.2000}}
  >{\raggedright\arraybackslash}p{(\linewidth - 8\tabcolsep) * \real{0.2000}}@{}}
\toprule\noalign{}
\begin{minipage}[b]{\linewidth}\raggedright
Client
\end{minipage} &
\begin{minipage}[b]{\linewidth}\raggedright
TLS
\end{minipage} &
\begin{minipage}[b]{\linewidth}\raggedright
Protocol
\end{minipage} &
\begin{minipage}[b]{\linewidth}\raggedright
Key's server got
\end{minipage} &
\begin{minipage}[b]{\linewidth}\raggedright
Stranger got
\end{minipage} \\
\midrule\noalign{}
\endhead
\bottomrule\noalign{}
\endlastfoot
Python http.client, requests, httpx &
system OpenSSL & HTTP/1.1, header & key
& nothing \\
Python & system OpenSSL & HTTP/1.1,
chunked body & key & nothing \\
Python & system OpenSSL & HTTP/1.1, gzip
body & placeholder & placeholder \\
Python & system OpenSSL & WebSocket
upgrade header & key & nothing \\
Python & system OpenSSL & WebSocket
message (masked) & placeholder &
placeholder \\
curl & system OpenSSL & HTTP/1.1 header;
Basic auth & key & nothing \\
curl & system OpenSSL, nghttp2 & HTTP/2
header (HPACK) & key & nothing \\
curl & system OpenSSL, nghttp2 & HTTP/2
body & key & the request, without key or
placeholder \\
curl & none & plain HTTP to a stranger &
-- & nothing \\
Ruby Net::HTTP & system OpenSSL &
HTTP/1.1, header & key & nothing \\
Node.js https & Node's OpenSSL &
HTTP/1.1, header & key & nothing \\
Go net/http & Go crypto/tls & HTTP/1.1,
header & placeholder & placeholder \\
\end{longtable}
}

Twelve of the 15 cells are protected.
The three that aren't (gzip bodies,
masked WebSocket messages, Go) never
release the key: the request to the
key's server fails, and the stranger
gets a placeholder. Not tested here:
Postgres and Redis over TLS, Java,
rustls, HTTP/3, and Windows.

\subsection*{5.2 Attack
replays}\label{attack-replays}
\addcontentsline{toc}{subsection}{5.2
Attack replays}

These are reconstructions of documented
behaviour, not live malware. The
Shai-Hulud replays send the
double-base64 haul to a webhook, folding
the worm's two channels (a public
repository and a webhook) into one. They
run locally against a synthetic key;
\texttt{eval/replays/README.md} has the
ethics note.

{\def\LTcaptype{none} % do not increment counter
\begin{longtable}[]{@{}
  >{\raggedright\arraybackslash}p{(\linewidth - 6\tabcolsep) * \real{0.2500}}
  >{\raggedright\arraybackslash}p{(\linewidth - 6\tabcolsep) * \real{0.2500}}
  >{\raggedright\arraybackslash}p{(\linewidth - 6\tabcolsep) * \real{0.2500}}
  >{\raggedright\arraybackslash}p{(\linewidth - 6\tabcolsep) * \real{0.2500}}@{}}
\toprule\noalign{}
\begin{minipage}[b]{\linewidth}\raggedright
Replay
\end{minipage} &
\begin{minipage}[b]{\linewidth}\raggedright
Attacker held
\end{minipage} &
\begin{minipage}[b]{\linewidth}\raggedright
Left the machine
\end{minipage} &
\begin{minipage}[b]{\linewidth}\raggedright
Stopped and attributed
\end{minipage} \\
\midrule\noalign{}
\endhead
\bottomrule\noalign{}
\endlastfoot
Nx ``s1ngularity'': env, \texttt{.env},
AI-CLI inventory, triple base64,
GitHub-style upload & placeholder &
nothing & yes \\
Shai-Hulud: env, double base64, HTTPS
webhook (Node) & placeholder & nothing &
yes \\
Shai-Hulud to a plain HTTP webhook &
placeholder (masked) & nothing & yes \\
\texttt{.env} hunter, posted in clear &
placeholder & nothing & yes \\
Agent (under a parent named
\texttt{claude}): print / base64 the key
& masked output & -- & -- (nothing
sent) \\
Agent:
\texttt{env\ \textbackslash{}\textbar{}\ curl}
to a paste site & masked output &
nothing & yes \\
Agent: post to a webhook & masked output
& nothing & yes \\
Agent: key in a DNS name & masked output
& -- & yes \\
Agent: write the key into a committed
file & placeholder & -- & -- (traced to
its run) \\
Use the key at its own server & -- & the
request, as you & no (out of scope) \\
\end{longtable}
}

Afterwards, a teammate on the same team
folder was refused every stopped
destination. The loot from the Nx,
Shai-Hulud, \texttt{.env} hunter and
commit replays traced back to the run
that produced it.

\subsection*{5.3
Overhead}\label{overhead}
\addcontentsline{toc}{subsection}{5.3
Overhead}

A Python client keeps one TLS connection
open to a local server. It sends 2,000
small requests after 50 warm-up ones,
then uploads 64 MiB in 64 KiB writes.
Figures are the median of 9 alternating
rounds. The baseline (``direct'') is
started the same way: detached, through
the runner, with no hook.

{\def\LTcaptype{none} % do not increment counter
\begin{longtable}[]{@{}
  >{\raggedright\arraybackslash}p{(\linewidth - 4\tabcolsep) * \real{0.3333}}
  >{\raggedright\arraybackslash}p{(\linewidth - 4\tabcolsep) * \real{0.3333}}
  >{\raggedright\arraybackslash}p{(\linewidth - 4\tabcolsep) * \real{0.3333}}@{}}
\toprule\noalign{}
\begin{minipage}[b]{\linewidth}\raggedright
\end{minipage} &
\begin{minipage}[b]{\linewidth}\raggedright
Per request
\end{minipage} &
\begin{minipage}[b]{\linewidth}\raggedright
Upload
\end{minipage} \\
\midrule\noalign{}
\endhead
\bottomrule\noalign{}
\endlastfoot
Direct & 61--63 µs & \textasciitilde940
MiB/s \\
Hook, nothing to swap & +2 µs & within
noise (−3\%) \\
Hook, key swapped in every request &
+262 µs & within noise (−13\% against a
±20\% spread) \\
\end{longtable}
}

The swap's cost is one local round trip
to the key service for the key; the
chain check is cached per connection.
Measuring forced two fixes:

\begin{itemize}
\tightlist
\item
  \textbf{Uploads first ran 80\% slower}
  under the hook, because 47 patterns
  were each searched separately, on both
  the plaintext and the ciphertext. The
  single-pass hashed scan fixed that.
\item
  \textbf{A swapped request first cost
  about 730 µs,} because the chain was
  verified every time and the parent
  chain was walked on every connection.
\end{itemize}

We have not yet compared
\texttt{op\ run} or a placeholder proxy,
and the planned week of ordinary
development under the hook has not been
run.

\subsection*{5.4 Adversarial
review}\label{adversarial-review}
\addcontentsline{toc}{subsection}{5.4
Adversarial review}

\texttt{docs/threats.md} tabulates 21
attacks by an attacker running as the
user. Each is classed as stopped,
worthless (only a placeholder leaves),
gets the key, or out of scope, with a
test or a reason.

The review found one real hole, now
fixed. The hook removed its token with
\texttt{unsetenv}, but the token stayed
in the process's first environment
block, readable at
\texttt{/proc/self/environ}. Code in the
program could read it and ask the key
service for the key. The test that
demonstrates this
(\texttt{attacks\_on\_placeholders\_get\_no\_key})
fails without the fix.

These attacks still get the key, by
design or as stated limits:

\begin{itemize}
\tightlist
\item
  \textbf{A token in a child without the
  hook.} A statically linked binary's
  environment keeps the token, and
  in-run code could read it there.
\item
  \textbf{In-process memory reads}
  during a swap.
\item
  \textbf{Use at the key's own server.}
\end{itemize}

No one outside the project has reviewed
the table yet.

\section*{6.
Discussion}\label{discussion}
\addcontentsline{toc}{section}{6.
Discussion}

\subsection*{6.1 What placeholders
change}\label{what-placeholders-change}
\addcontentsline{toc}{subsection}{6.1
What placeholders change}

The common theft paths become harmless:
reading the environment or
\texttt{.env}, printing, uploading,
pasting into an agent transcript,
committing. Each yields a placeholder
that dies with its run, and every send
of one through a covered path is stopped
and attributed.

What remains open is narrower and
different in kind. The attacker must
either use the key through its own
legitimate server, which is visible in
that provider's logs, or write code that
targets the hook specifically.

\subsection*{6.2 Why the hook, and not a
proxy}\label{why-the-hook-and-not-a-proxy}
\addcontentsline{toc}{subsection}{6.2
Why the hook, and not a proxy}

A proxy sees traffic only if the program
routes it there and trusts the proxy's
CA. Code in the program can undo both.
It also understands only the protocols
the proxy parses.

The in-process hook works on any
protocol over the TLS libraries it
covers, with no routing and no new CA.
It decides on the certificate the server
actually presented, not on what the
program asked for.

The cost is that the key is in the
program's memory during the swap. Moving
the swap out of the process, with an
eBPF uprobe on \texttt{SSL\_write} as
Kloak does in Kubernetes {[}14{]}, would
close that gap. It needs privileges a
developer's desktop tool shouldn't ask
for.

\subsection*{6.3 Why tripwires have no
false
positives}\label{why-tripwires-have-no-false-positives}
\addcontentsline{toc}{subsection}{6.3
Why tripwires have no false positives}

A placeholder is never a valid
credential anywhere. Nothing legitimate
can send it to a server other than its
key's. A refused write is therefore
always a real attempt.

The exceptions are configuration errors.
A wrong host on the key makes legitimate
requests fail closed rather than leak. A
teammate's warning about a service the
developer needs blocks it; warnings name
their signer.

\subsection*{6.4
Limitations}\label{limitations}
\addcontentsline{toc}{subsection}{6.4
Limitations}

\begin{itemize}
\tightlist
\item
  \textbf{Coverage.} Coverage is the TLS
  libraries hooked: dynamically linked
  OpenSSL and Node. Windows (schannel)
  is not built.
\item
  \textbf{Placeholders that aren't
  seen.} Placeholders inside compression
  or other encodings aren't recognised,
  and neither are header blocks split
  across HTTP/2 CONTINUATION frames.
\item
  \textbf{Agent detection} is by process
  name.
\item
  \textbf{Evaluation.} The replays are
  reconstructions, run on one machine
  rather than in disposable VMs. The
  false-alarm week hasn't happened.
\end{itemize}

\section*{7. Related
work}\label{related-work}
\addcontentsline{toc}{section}{7.
Related work}

\begin{itemize}
\tightlist
\item
  \textbf{Run wrappers.} 1Password's
  \texttt{op\ run} {[}4{]},
  \texttt{doppler\ run} {[}5{]} and
  \texttt{infisical\ run} {[}6{]} inject
  keys into a child's environment.
  \texttt{op\ run} also masks secrets
  the child prints. The child holds the
  real key.
\item
  \textbf{Placeholder proxies.} Docker's
  agent sandboxes {[}7{]}, Infisical's
  Agent Vault {[}8{]}, iron-proxy
  {[}9{]}, fnox's proxy {[}10{]} and
  GhostKey {[}11{]} give the workload a
  placeholder and swap in the real key
  at an HTTP(S) proxy, which the
  workload must route through and whose
  CA it must trust. fnox notes that a
  same-user process can bypass the
  proxy's environment settings.
\item
  \textbf{Sandbox credential brokers.}
  Bashkit's placeholder mode {[}12{]}
  swaps a placeholder in outbound HTTP
  headers for scripts it runs, and
  OpenSandbox's credential vault
  {[}13{]} keeps the real key in an
  egress sidecar outside the sandbox. A
  design proposed for Sandlock {[}15{]}
  keeps the secret in a supervisor
  process and gives the child a
  per-session phantom token. All of
  these protect a sandboxed guest; ours
  protects programs a developer runs on
  their own machine.
\item
  \textbf{Kernel uprobes.} Kloak
  {[}14{]} swaps placeholders for
  Kubernetes workloads with an eBPF
  uprobe on \texttt{SSL\_write},
  restricted to allowed destinations. It
  is the closest mechanism to ours, at a
  different layer and for a different
  user.
\item
  \textbf{eBPF TLS tracing.} Tools such
  as eCapture {[}16{]} read
  \texttt{SSL\_write} plaintext through
  uprobes without changing it.
\item
  \textbf{Honeytokens.} Decoy
  credentials that alarm when used
  {[}18{]} are the ancestor of our
  tripwires, and per-site unique decoys
  can show where a leak happened
  {[}19{]}. Honeywords {[}20{]} plant
  decoy passwords beside the real one.
  Our placeholders differ in that they
  are the credential the program
  actually uses, fresh per run.
\item
  \textbf{Artificial immune systems.}
  Forrest, Hofmeyr and others applied
  self/non-self discrimination to
  intrusion detection {[}21, 22{]}.
\end{itemize}

\section*{8. Ethics}\label{ethics}
\addcontentsline{toc}{section}{8.
Ethics}

The replays reimplement only the
collection and exfiltration steps
described in public reports. They have
no propagation, persistence or
destructive behaviour. They run against
local servers with synthetic keys, and
\texttt{.test} names never resolve
outside the machine.

\section*{9. Artifact}\label{artifact}
\addcontentsline{toc}{section}{9.
Artifact}

Run \texttt{eval/run.sh} from a
checkout. It needs Rust, Node.js, Python
3, openssl, curl and setsid; Ruby, Go,
\texttt{requests} and \texttt{httpx} add
their rows when present.

It writes
\texttt{eval/results/coverage.md},
\texttt{replays.md} and
\texttt{benchmarks.md}, with JSON beside
each. CI
(\texttt{.github/workflows/check.yml})
reruns the coverage matrix and the
replays on every push, and fails if a
claimed cell isn't protected, a stranger
ever gets a key, or a replayed send
isn't stopped and attributed.

\section*{Acknowledgments}\label{acknowledgments}
\addcontentsline{toc}{section}{Acknowledgments}

Claude (Anthropic), an AI model, was
used extensively to write the software,
run the evaluation and draft this paper.
The author reviewed the text, checked
every number against
\texttt{eval/results/} and every
reference against its source, and takes
responsibility for the contents.

\section*{References}\label{references}
\addcontentsline{toc}{section}{References}

{[}1{]} Wiz Research. s1ngularity:
supply chain attack leaks secrets on
GitHub: everything you need to know.
August 2025.
\url{https://www.wiz.io/blog/s1ngularity-supply-chain-attack}

{[}2{]} Wiz Research. s1ngularity's
aftermath: AI, TTPs, and impact in the
Nx supply chain attack. 2025.
\url{https://www.wiz.io/blog/s1ngularitys-aftermath}

{[}3{]} Wiz Research. Shai-Hulud:
ongoing package supply chain worm
delivering data-stealing malware.
September 2025.
\url{https://www.wiz.io/blog/shai-hulud-npm-supply-chain-attack}

{[}4{]} 1Password. op run (CLI
reference).
\url{https://www.1password.dev/cli/reference/commands/run}

{[}5{]} Doppler. CLI guide.
\url{https://docs.doppler.com/docs/cli}

{[}6{]} Infisical. infisical run (CLI
reference).
\url{https://infisical.com/docs/cli/commands/run}

{[}7{]} Docker. Docker Sandboxes:
credentials.
\url{https://docs.docker.com/ai/sandboxes/security/credentials/}

{[}8{]} Infisical. Agent Vault.
\url{https://github.com/Infisical/agent-vault}

{[}9{]} ironsh. iron-proxy.
\url{https://github.com/ironsh/iron-proxy}

{[}10{]} fnox. Proxy guide.
\url{https://fnox.jdx.dev/guide/proxy.html}

{[}11{]} GhostKey.
\url{https://github.com/jhaji2911/GhostKey}

{[}12{]} Bashkit. Credential injection.
\url{https://docs.rs/crate/bashkit/latest/source/docs/credential-injection.md}

{[}13{]} OpenSandbox. Credential vault.
\url{https://github.com/opensandbox-group/opensandbox}

{[}14{]} spinningfactory. Kloak.
\url{https://github.com/spinningfactory/kloak}

{[}15{]} multikernel. RFC: credential
injection via the proxy (Sandlock issue
66).
\url{https://github.com/multikernel/sandlock/issues/66}

{[}16{]} gojue. eCapture.
\url{https://github.com/gojue/ecapture}

{[}17{]} R. Peon and H. Ruellan. HPACK:
Header Compression for HTTP/2. RFC 7541,
IETF, 2015.

{[}18{]} L. Spitzner. Honeytokens: the
other honeypot. SecurityFocus, 2003.

{[}19{]} J. DeBlasio, S. Savage, G. M.
Voelker and A. C. Snoeren. Tripwire:
inferring internet site compromise. In
Proc. ACM Internet Measurement
Conference (IMC), 2017.

{[}20{]} A. Juels and R. L. Rivest.
Honeywords: making password-cracking
detectable. In Proc. ACM CCS, 2013.

{[}21{]} S. Forrest, S. A. Hofmeyr, A.
Somayaji and T. A. Longstaff. A sense of
self for Unix processes. In Proc. IEEE
Symposium on Security and Privacy, 1996.

{[}22{]} S. A. Hofmeyr and S. Forrest.
Architecture for an artificial immune
system. Evolutionary Computation 8(4),
2000.

\end{document}
